\begin{center}
\section*{A}
\end{center}
\addcontentsline{toc}{section}{A}
\term{Agile}
Agile è un approccio di lavoro basato su iterazioni brevi e feedback. Utilizzato dal metodo Scrum, il lavoro viene suddiviso in Sprint di una o due settimane, ciascun Sprint inizia con lo Sprint Planning, dove vengono definiti obiettivi, e si conclude con lo Sprint Retrospective, dove il team valuta cosa è andato bene e cosa è necessario migliorare. Durante lo Sprint vengono svolte riunioni giornaliere per monitorare lo stato di avanzamento.
\term{Amministratore}
È la figura che si occupa della configurazione e gestione dell'infrastruttura IT di supporto a un progetto.
\term{Analisi dei Requisiti}
L'analisi dei requisiti è il processo di raccolta, documentazione e analisi delle esigenze e delle funzionalità di un sistema software. Questo processo coinvolge la comprensione approfondita dei requisiti dell'utente, delle necessità del business e delle specifiche tecniche per sviluppare un sistema che soddisfi le aspettative degli stakeholders. Gli obiettivi principali dell'analisi dei requisiti includono la definizione chiara e completa delle funzionalità del sistema, la specifica dei vincoli e delle restrizioni, e l'identificazione dei requisiti non funzionali come quelli relativi a prestazioni, sicurezza e usabilità.
\term{Analisi dei Rischi}
L'analisi dei rischi è il processo di identificazione, valutazione e gestione dei potenziali problemi e delle minacce che potrebbero influenzare il successo di un progetto o di un'attività. Questo processo coinvolge la valutazione dei rischi potenziali, la determinazione della loro probabilità di accadere e dell'impatto che potrebbero avere sul progetto, nonché lo sviluppo di strategie per mitigare o gestire tali rischi. L'obiettivo dell'analisi dei rischi è identificare proattivamente le potenziali problematiche e adottare misure preventive o correttive per ridurre al minimo gli impatti negativi sul progetto.
\term{Analista}
Figura che si occupa di identificare e chiarire i requisiti, interpretando le esigenze degli utilizzatori finali così da assicurare una corretta definizione delle funzionalità.
\term{API}
Acronimo di Application Programming Interface, è un'interfaccia utile a permettere o facilitare la comunicazione tra software o parti di un singolo software.
\term{Attore}
Entità esterna e non controllabile dal progetto ma che va ad intergaire con esso.
